\chapter{Requirements Analysis and Scrum Methodology}

\section*{Introduction to Chapter 3}
\addcontentsline{toc}{section}{Introduction to Chapter 3}

This chapter presents the requirements analysis and the Scrum methodology adopted for the project. It first identifies the three actors who will interact with the platform: the Administrator, the Maintenance Engineer, and the Drone Operator. It then describes their needs, followed by the functional and non-functional requirements. The chapter also defines the project scope and justifies the monitoring parameters selected. The Scrum roles are assigned, the product backlog is presented through epics and user stories, the definition of done is established, and the chapter concludes with the sprint planning for the five sprints.

\section{Actors Identification}

The following table presents the three actors of the system:

\begin{table}[H]
\centering
\caption{Actors of the system}
\begin{tabular}{|p{3cm}|p{4cm}|p{6cm}|}
\hline
\textbf{Actor} & \textbf{Description} & \textbf{What they need to do} \\
\hline
\textbf{Administrator} & Person who manages the platform and user accounts & Create user accounts, disable user accounts, view all system activity, configure settings \\
\hline
\textbf{Maintenance Engineer} & Person who fixes the drone and analyzes health data & View telemetry graphs, see AI anomaly detection results, check historical flights, read anomaly explanations, plan repairs \\
\hline
\textbf{Drone Operator} & Person who flies the drone & Monitor motor temperature and current during flight, receive real-time alerts, land safely if a problem is detected, inform the maintenance engineer of sudden failures \\
\hline
\end{tabular}
\end{table}

\section{Needs Analysis}

\subsection{Administrator Needs}

The Administrator needs to manage user accounts. When a new Maintenance Engineer or Drone Operator joins AVIONAV, the Administrator must be able to create an account for them. When someone leaves the company, the Administrator must be able to disable their account. The Administrator also needs to view system activity and configure platform settings.

\subsection{Maintenance Engineer Needs}

The Maintenance Engineer needs to monitor the health of the UAV motor. They need to view real-time and historical telemetry graphs (current and temperature). They need to see when the AI detects an anomaly and understand why the anomaly was detected (explainability). They need to plan maintenance actions (such as bearing replacement) based on the AI predictions.

\subsection{Drone Operator Needs}

The Drone Operator needs to receive real-time alerts during flight. They need to know if the motor is becoming too hot or pulling too much current. They need to land the drone safely if an alert is triggered. If a sudden failure occurs (like a wire snapping), they need to inform the Maintenance Engineer.

\section{Scope Definition and Justification}

The initial technical specification proposed monitoring five avionics-related parameters: electrical voltage levels, electrical current measurements, temperature sensor data, error rate counters, and communication message frequency. However, after discussion with the industrial supervisor, the project scope was focused specifically on the UAV propulsion system (brushless DC motor and electronic speed controller). The motor was identified as the most critical component of a fixed-wing UAV, as its failure leads to immediate loss of propulsion and potential crash.

Among the possible motor failure modes (broken propeller, burnt ESC, disconnected wire, demagnetised magnets, bearing wear), bearing wear was selected as the focus because it represents a progressive mechanical failure that generates time-series data patterns suitable for machine learning detection. Bearing wear occurs gradually over hours or days, producing slowly rising current and temperature signatures that can be detected before catastrophic failure occurs. Therefore, the monitoring scope was reduced to two parameters: motor current (Amperes) and motor temperature (degrees Celsius). These two parameters, combined through the physical relationship described by thermal dynamics (Ploss = I$^2$ × R), provide sufficient information to detect bearing wear degradation patterns.

\section{Functional and Non-Functional Requirements}

\subsection{Functional Requirements}

\begin{table}[H]
\centering
\caption{Functional requirements}
\begin{tabular}{|p{1.5cm}|p{10cm}|}
\hline
\textbf{ID} & \textbf{Requirement} \\
\hline
FR-01 & The system shall allow users to log in with an email and password. \\
\hline
FR-02 & The system shall allow the Administrator to create, view, and disable user accounts. \\
\hline
FR-03 & The system shall generate simulated telemetry data for motor current (Amps) and motor temperature (°C). \\
\hline
FR-04 & The system shall store telemetry data in a database (SQLite). \\
\hline
FR-05 & The system shall train an Isolation Forest machine learning model on healthy flight data. \\
\hline
FR-06 & The system shall detect anomalies in real-time telemetry streams using the trained model. \\
\hline
FR-07 & The system shall display current and temperature graphs on a dashboard. \\
\hline
FR-08 & The system shall send an alert when an anomaly is detected. \\
\hline
FR-09 & The system shall provide an explanation of why an anomaly was detected. \\
\hline
FR-10 & The system shall allow the Maintenance Engineer to view historical flight data. \\
\hline
FR-11 & The system shall distinguish between three user roles (Administrator, Maintenance Engineer, Drone Operator). \\
\hline
FR-12 & The system shall allow the Drone Operator to view only alerts and basic status (not full telemetry history). \\
\hline
\end{tabular}
\end{table}

\subsection{Non-Functional Requirements}

\begin{table}[H]
\centering
\caption{Non-functional requirements}
\begin{tabular}{|p{1.5cm}|p{10cm}|}
\hline
\textbf{ID} & \textbf{Requirement} \\
\hline
NFR-01 & The system shall respond to user actions within 2 seconds. \\
\hline
NFR-02 & The system shall require authentication for all pages except the login page. \\
\hline
NFR-03 & The system shall be accessible via a web browser (Chrome, Firefox, Edge). \\
\hline
NFR-04 & The system shall run on a standard laptop without special hardware. \\
\hline
NFR-05 & The system shall store data persistently using SQLite so that data survives restarts. \\
\hline
NFR-06 & The system shall minimize false negatives (missing a real failure) to ensure drone safety. \\
\hline
\end{tabular}
\end{table}

\section{Scrum Roles}

In accordance with the Scrum framework, the project is structured around three main roles:

\begin{itemize}
    \item \textbf{Product Owner:} Assured by the Academic Supervisor (Mr. Aymen Charrada). Responsible for defining the project needs, validating the functional increments, and prioritizing the objectives.
    \item \textbf{Scrum Master:} Assured by the student (Mootez Bourguiba). Responsible for organizing the work, applying Scrum principles, and ensuring the smooth execution of Sprints.
    \item \textbf{Development Team:} Assured by the student (Mootez Bourguiba). Responsible for the technical conception, Python development, AI training, and testing of the platform.
\end{itemize}

\section{Product Backlog: Epics and User Stories}

\subsection{Defined Epics}

\begin{table}[H]
\centering
\caption{Epics of the project}
\begin{tabular}{|p{2cm}|p{8cm}|p{3cm}|}
\hline
\textbf{Epic ID} & \textbf{Description} & \textbf{Associated Sprint} \\
\hline
Epic 1 & Scope Definition and Justification & Sprint 0 \\
\hline
Epic 2 & UAV Context and Configuration & Sprint 0 \\
\hline
Epic 3 & System Architecture Design & Sprint 0 \\
\hline
Epic 4 & User Authentication and Management & Sprint 1 \\
\hline
Epic 5 & Database and Data Model & Sprint 1 \\
\hline
Epic 6 & Telemetry Data Pipeline & Sprint 2 \\
\hline
Epic 7 & Anomaly Detection Engine & Sprint 3 \\
\hline
Epic 8 & REST API & Sprint 4 \\
\hline
Epic 9 & Streamlit Dashboard & Sprint 4 \\
\hline
Epic 10 & System Integration and Testing & Sprint 5 \\
\hline
\end{tabular}
\end{table}

\subsection{User Stories}

\begin{table}[H]
\centering
\caption{Product backlog user stories}
\begin{tabular}{|p{1.5cm}|p{2.5cm}|p{5cm}|p{1.5cm}|p{1.5cm}|}
\hline
\textbf{ID} & \textbf{Epic} & \textbf{User Story} & \textbf{Priority} & \textbf{Story Points} \\
\hline
US-00 & Epic 1 & Document scope decisions (reduce from 5 to 2 parameters based on industrial supervisor guidance) & High & 2 \\
\hline
US-01 & Epic 2 & Provide UAV context (basic configuration selection for motor monitoring) & High & 3 \\
\hline
US-02 & Epic 3 & Design software architecture and database schema & High & 5 \\
\hline
US-03 & Epic 4 & As a user, I want to log in with my email and password so that I can access the platform. & High & 5 \\
\hline
US-04 & Epic 4 & As an Administrator, I want to create new user accounts for Maintenance Engineers and Drone Operators. & High & 5 \\
\hline
US-05 & Epic 4 & As an Administrator, I want to disable user accounts when someone leaves AVIONAV. & High & 3 \\
\hline
US-06 & Epic 5 & As a system, I want to store user accounts in database & High & 3 \\
\hline
US-07 & Epic 5 & As a system, I want to store flight sessions in database & High & 3 \\
\hline
US-08 & Epic 5 & As a system, I want to store telemetry readings in database & High & 3 \\
\hline
US-09 & Epic 5 & As a system, I want to store anomaly events in database & High & 3 \\
\hline
US-10 & Epic 6 & As a Maintenance Engineer, I want the system to ingest telemetry data from CSV files & Medium & 5 \\
\hline
US-11 & Epic 6 & As a system, I want to preprocess telemetry data (cleaning, normalization) & Medium & 3 \\
\hline
US-12 & Epic 7 & As a Maintenance Engineer, I want the Isolation Forest model to be trained on healthy flight data only. & High & 5 \\
\hline
US-13 & Epic 7 & As a Maintenance Engineer, I want the system to detect anomalies (gradual current and temperature drift) on new flight data. & High & 5 \\
\hline
US-14 & Epic 8 & As a client, I want to access flight data via REST API & High & 5 \\
\hline
US-15 & Epic 9 & As a Maintenance Engineer, I want to view dashboard with telemetry charts & High & 8 \\
\hline
US-16 & Epic 9 & As a Maintenance Engineer, I want to receive visual alerts for anomalies & High & 5 \\
\hline
\end{tabular}
\end{table}

\subsection{Acceptance Criteria}

\begin{table}[H]
\centering
\caption{Acceptance criteria for user stories}
\begin{tabular}{|p{1.5cm}|p{10cm}|}
\hline
\textbf{US ID} & \textbf{Acceptance Criteria} \\
\hline
US-00 & Scope document is written and approved by industrial supervisor. Justification for parameter reduction is clearly documented. \\
\hline
US-01 & UAV configuration is selected (motor type, power rating). Basic specifications are documented. \\
\hline
US-02 & Software architecture diagram is created. Database schema (ER diagram) is designed. Technology stack is justified. \\
\hline
US-03 & Login form accepts email and password. JWT token is generated on successful login. Error handling for invalid credentials. \\
\hline
US-04 & Admin can create user with email, password, and role. Password validation (min 8 characters). Success/error feedback. \\
\hline
US-05 & Admin can disable user account. Confirmation dialog. User cannot login after disable. \\
\hline
US-06 & User table exists with required fields. SQLAlchemy ORM model is defined. Alembic migration is created and tested. \\
\hline
US-07 & FlightSession table exists with foreign key to User. Alembic migration is created and tested. \\
\hline
US-08 & TelemetryReading table exists with foreign key to FlightSession. Alembic migration is created and tested. \\
\hline
US-09 & AnomalyEvent table exists with foreign key to FlightSession. Alembic migration is created and tested. \\
\hline
US-10 & CSV file is parsed correctly. Data validation (current and temperature ranges). Data is inserted into database. \\
\hline
US-11 & Null values are removed. Outliers are filtered. Data is normalized. Preprocessed data is stored or passed to next stage. \\
\hline
US-12 & Isolation Forest is trained on healthy flights only. Model is saved to disk. Training metrics are recorded. \\
\hline
US-13 & Model is loaded and used for prediction. Anomaly scores are calculated. Anomalies above threshold are stored in database. \\
\hline
US-14 & GET /api/flights returns list of flights. GET /api/flights/{id} returns specific flight. JWT authentication is required. \\
\hline
US-15 & Dashboard displays current and temperature line charts. Flight selector is available. Real-time updates work. \\
\hline
US-16 & Color-coded alerts (green/yellow/red) are displayed. Alert details panel shows severity. \\
\hline
\end{tabular}
\end{table}

\section{Definition of Done}

A User Story is considered \textbf{Done} when the following criteria are met:

\begin{itemize}
    \item All acceptance criteria are satisfied
    \item Code is implemented and follows PEP 8 style guidelines
    \item Unit tests are written and pass with at least 80\% coverage
    \item Code is reviewed and approved
    \item API documentation is updated (if applicable)
    \item Integration with existing system is verified
    \item No critical bugs remain
    \item Security requirements are verified (no vulnerabilities introduced)
    \item Performance meets non-functional requirements
\end{itemize}

A Sprint is considered \textbf{Done} when:

\begin{itemize}
    \item All User Stories in the Sprint are Done according to the Definition of Done
    \item The Sprint Goal is achieved
    \item A functional increment is delivered and can be demonstrated
    \item Sprint chapter is drafted for the report
    \item UML diagrams (Use Case, Class, Sequence) are created for the Sprint
    \item Sprint Review is conducted and demo is presented
\end{itemize}

\section{Requirements Traceability Matrix}

\begin{table}[H]
\centering
\caption{Requirements traceability matrix}
\begin{tabular}{|p{2cm}|p{2cm}|p{6cm}|}
\hline
\textbf{Functional Req.} & \textbf{User Story} & \textbf{Epic} \\
\hline
FR-01 & US-03 & Epic 4 \\
\hline
FR-02 & US-04, US-05 & Epic 4 \\
\hline
FR-03 & US-10 & Epic 6 \\
\hline
FR-04 & US-06, US-07, US-08, US-09 & Epic 5 \\
\hline
FR-05 & US-12 & Epic 7 \\
\hline
FR-06 & US-13 & Epic 7 \\
\hline
FR-07 & US-15 & Epic 9 \\
\hline
FR-08 & US-16 & Epic 9 \\
\hline
FR-09 & US-13 (explainability via feature importance) & Epic 7 \\
\hline
FR-10 & US-14, US-15 & Epic 8, Epic 9 \\
\hline
FR-11 & US-04, US-05 & Epic 4 \\
\hline
FR-12 & US-15, US-16 (role-based views) & Epic 9 \\
\hline
\end{tabular}
\end{table}

\section{Sprint Planning}

The project timeline is divided into five implementation sprints plus a final integration sprint:

\begin{table}[H]
\centering
\caption{Sprint planning}
\begin{tabular}{|p{2cm}|p{3cm}|p{8cm}|}
\hline
\textbf{Sprint} & \textbf{Duration} & \textbf{Objective / Deliverable} \\
\hline
Sprint 0 & Weeks 1–2 & Initialization: Scope definition, UAV context, software architecture, database schema design \\
\hline
Sprint 1 & Weeks 3–4 & Authentication and Database: User authentication, user management, database setup with all tables \\
\hline
Sprint 2 & Weeks 5–6 & Data Pipeline: CSV telemetry ingestion, data preprocessing, simulated dataset generation \\
\hline
Sprint 3 & Weeks 7–8 & Anomaly Detection: Isolation Forest model training, anomaly detection inference, validation \\
\hline
Sprint 4 & Weeks 9–11 & API and Dashboard: REST API endpoints, Streamlit dashboard, real-time visualization, alerts \\
\hline
Sprint 5 & Weeks 12–14 & Integration and Report: End-to-end testing, bug fixes, report finalization, presentation preparation \\
\hline
\end{tabular}
\end{table}

\section*{Conclusion to Chapter 3}
\addcontentsline{toc}{section}{Conclusion to Chapter 3}

This chapter has presented the requirements analysis for the Intelligent Avionics Health Monitoring Platform. Three actors were identified: Administrator, Maintenance Engineer, and Drone Operator. Their needs were analyzed, and functional and non-functional requirements were defined. The project scope was justified, explaining the focus on motor current and temperature monitoring for bearing wear detection. The Scrum roles were assigned, and the product backlog was presented through ten epics and sixteen user stories with acceptance criteria and story points. The definition of done was established, and the requirements traceability matrix was provided. Finally, the sprint planning was established, dividing the project into five implementation sprints plus a final integration sprint. The next chapter will present Sprint 0: Initialization and Architecture.
